\documentclass[11pt,a4paper]{article}

\usepackage[margin=1in]{geometry}
\usepackage{listings}
\usepackage{xcolor}
\usepackage{tikz}
\usetikzlibrary{shapes.geometric,arrows.meta,positioning,fit,backgrounds}
\usepackage{booktabs}
\usepackage{hyperref}
\usepackage{caption}
\usepackage{amsmath}
\usepackage{enumitem}
\usepackage{fancyhdr}
\usepackage{titling}

\hypersetup{colorlinks=true, linkcolor=blue, urlcolor=blue, citecolor=blue}

\definecolor{codebg}{rgb}{0.96,0.96,0.96}
\definecolor{codekw}{rgb}{0.0,0.0,0.6}
\definecolor{codecomment}{rgb}{0.3,0.5,0.3}
\definecolor{codestring}{rgb}{0.6,0.0,0.0}

\lstdefinestyle{base}{
  backgroundcolor=\color{codebg},
  basicstyle=\ttfamily\small,
  commentstyle=\color{codecomment}\itshape,
  keywordstyle=\color{codekw}\bfseries,
  stringstyle=\color{codestring},
  numbers=left,
  numberstyle=\tiny\color{gray},
  numbersep=8pt,
  frame=single,
  rulecolor=\color{gray!40},
  breaklines=true,
  breakatwhitespace=false,
  showstringspaces=false,
  tabsize=2,
  columns=fullflexible,
  keepspaces=true
}
\lstset{style=base}

\lstdefinelanguage{pseudo}{
  morekeywords={if,else,for,in,repeat,return,raise,function,not,every,delete,goto,while},
  morecomment=[l]{//},
  sensitive=true
}
\lstdefinestyle{pseudostyle}{
  language=pseudo,
  backgroundcolor=\color{codebg},
  basicstyle=\ttfamily\small,
  commentstyle=\color{codecomment}\itshape,
  keywordstyle=\color{codekw}\bfseries,
  frame=single,
  rulecolor=\color{gray!40},
  breaklines=true,
  breakatwhitespace=false,
  showstringspaces=false,
  tabsize=2,
  columns=fullflexible,
  keepspaces=true
}

\pagestyle{fancy}
\fancyhf{}
\fancyhead[L]{Branch and Loop Handling in the APARA C Compiler}
\fancyhead[R]{\thepage}
\renewcommand{\headrulewidth}{0.4pt}

\setlength{\droptitle}{-1.5cm}
\title{\textbf{Branch and Loop Handling in the APARA C Compiler}\\[2mm]
\large How C control flow is broken down into APARA mcode instructions}
\author{APARA C Compiler Project}
\date{\today}

\begin{document}
\maketitle
\thispagestyle{fancy}
\vspace{-8mm}

% ============================================================================
\section{Introduction and Strategy}
% ============================================================================

This report explains how the compiler lowers every C control-flow construct ---
\texttt{if}/\texttt{else}, \texttt{while}, \texttt{do-while}, \texttt{for}, \texttt{switch},
\texttt{break}/\texttt{continue} --- down to APARA mcode, across the same two compiler
stages used for everything else: \texttt{ir\_gen.py} (AST $\to$ three-address IR) and
\texttt{codegen.py} (IR $\to$ mcode).

\textbf{The single governing strategy} is that the IR has no nested structure at all ---
no if/while/for nodes survive into it. Every construct is lowered, once and for all in
\texttt{ir\_gen.py}, into exactly three IR primitives: \texttt{IRLabel} (a named point in
the flat instruction stream), \texttt{IRJump} (an unconditional jump to a label), and
\texttt{IRCondJump} (jump to a label if a comparison holds). By the time \texttt{codegen.py}
sees the program, an \texttt{if} statement and a \texttt{while} loop are built from
\emph{exactly the same two ingredients} --- only the arrangement of labels and jumps
differs. This is what makes a single, uniform mcode-level treatment (Section~\ref{sec:mcode})
possible: the bundler and the assembler never need to know what kind of C statement
produced a given branch.

% ============================================================================
\section{Lowering Each Construct: Algorithms and Flow}
% ============================================================================

Each construct below is handled by its own \texttt{visit\_Xxx} method in
\texttt{ir\_gen.py}, but all of them follow the same pattern: allocate a small, fixed set of
labels up front, then emit \texttt{IRLabel}\allowbreak/\texttt{IRCondJump}\allowbreak/%
\texttt{IRJump} around a recursive call that lowers the nested body.

\subsection*{\texttt{if} / \texttt{if-else}}
\begin{lstlisting}[style=pseudostyle, caption={\texttt{visit\_If}, ir\_gen.py:724--735}]
function LOWER-IF(cond, ifTrueBody, ifFalseBody):
    t_lbl, e_lbl <- new labels "if_t", "if_e"
    if ifFalseBody exists:
        f_lbl <- new label "if_f"
        EMIT-COND(cond, true->t_lbl, false->f_lbl)
        LABEL t_lbl;  LOWER(ifTrueBody);  JUMP e_lbl
        LABEL f_lbl;  LOWER(ifFalseBody)
    else:
        EMIT-COND(cond, true->t_lbl, false->e_lbl)
        LABEL t_lbl;  LOWER(ifTrueBody)
    LABEL e_lbl
\end{lstlisting}

\begin{figure}[h]
\centering
\begin{tikzpicture}[
  node distance=8mm and 12mm,
  condshape/.style={diamond, draw, thick, aspect=2.2, minimum width=2.6cm, align=center, fill=blue!8},
  box/.style={rectangle, draw, thick, rounded corners, minimum width=2.6cm, minimum height=0.8cm, align=center, fill=orange!10},
  lbl/.style={rectangle, draw=none, fill=none, font=\scriptsize\ttfamily},
  arr/.style={-{Triangle[length=2.5mm,width=2mm]}, thick}
]
  \node[condshape] (cond) {cond};
  \node[box, below left=of cond] (true) {if\_t: \\ ifTrue body};
  \node[box, below right=of cond] (false) {if\_f: \\ ifFalse body};
  \node[box, below=18mm of cond] (end) {if\_e};
  \draw[arr] (cond) -- node[lbl, left]{true} (true);
  \draw[arr] (cond) -- node[lbl, right]{false} (false);
  \draw[arr] (true) -- (end);
  \draw[arr] (false) -- (end);
\end{tikzpicture}
\caption{\texttt{if}/\texttt{else} flow. The \texttt{if\_f} branch and its incoming edge
disappear entirely when there is no \texttt{else}.}
\end{figure}

\subsection*{\texttt{while}}
\begin{lstlisting}[style=pseudostyle, caption={\texttt{visit\_While}, ir\_gen.py:737--746}]
function LOWER-WHILE(cond, body):
    cond_lbl, body_lbl, end_lbl <- new labels "wc", "wb", "we"
    push (break->end_lbl, continue->cond_lbl)
    LABEL cond_lbl
    EMIT-COND(cond, true->body_lbl, false->end_lbl)
    LABEL body_lbl;  LOWER(body);  JUMP cond_lbl
    LABEL end_lbl
    pop break/continue targets
\end{lstlisting}

\begin{figure}[h]
\centering
\begin{tikzpicture}[
  node distance=10mm and 16mm,
  condshape/.style={diamond, draw, thick, aspect=2.2, minimum width=2.6cm, align=center, fill=blue!8},
  box/.style={rectangle, draw, thick, rounded corners, minimum width=2.8cm, minimum height=0.8cm, align=center, fill=orange!10},
  arr/.style={-{Triangle[length=2.5mm,width=2mm]}, thick}
]
  \node[condshape] (cond) {wc: cond};
  \node[box, right=of cond] (body) {wb: body};
  \node[box, below=of cond] (end) {we};
  \draw[arr] (cond) -- node[above]{true} (body);
  \draw[arr] (cond) -- node[left]{false} (end);
  \draw[arr] (body.south) to[out=-90,in=0] (end.east |- body.south) -- ++(0,-10mm) -| ([xshift=-1mm]cond.south);
  \draw[arr] (body) to[out=180,in=0] (cond);
\end{tikzpicture}
\caption{\texttt{while} flow: the body always jumps back to \texttt{wc}, which is the
only place the condition is (re-)tested. \texttt{break} jumps straight to \texttt{we};
\texttt{continue} jumps straight to \texttt{wc}.}
\end{figure}

\subsection*{\texttt{do-while}}
Structurally the same three labels as \texttt{while}, but the body label comes
\emph{first} and the condition is tested only after the first execution of the body ---
the condition check still lives at a single point, just after the body instead of before
it (\texttt{visit\_DoWhile}, ir\_gen.py:748--756):
\begin{lstlisting}[style=pseudostyle]
function LOWER-DO-WHILE(body, cond):
    body_lbl, cond_lbl, end_lbl <- new labels "dw_b", "dw_c", "dw_e"
    LABEL body_lbl;  LOWER(body)
    LABEL cond_lbl
    EMIT-COND(cond, true->body_lbl, false->end_lbl)
    LABEL end_lbl
\end{lstlisting}

\subsection*{\texttt{for}}
A \texttt{for} loop is a \texttt{while} loop with two extra, fixed slots: the
initializer runs once before the condition label, and the increment expression gets its
\emph{own} label so that \texttt{continue} can jump there without re-running the
initializer (\texttt{visit\_For}, ir\_gen.py:758--774):
\begin{lstlisting}[style=pseudostyle]
function LOWER-FOR(init, cond, incr, body):
    cond_lbl, body_lbl, incr_lbl, end_lbl <- new labels "fc", "fb", "fi", "fe"
    push (break->end_lbl, continue->incr_lbl)
    LOWER(init)                                    // runs exactly once
    LABEL cond_lbl
    if cond exists:  EMIT-COND(cond, true->body_lbl, false->end_lbl)
    else:             JUMP body_lbl                // no condition -> infinite loop
    LABEL body_lbl;  LOWER(body)
    LABEL incr_lbl;  LOWER(incr);  JUMP cond_lbl
    LABEL end_lbl
    pop break/continue targets
\end{lstlisting}

\subsection*{\texttt{switch}}
APARA mcode has no jump-table instruction, so \texttt{switch} lowers to a linear chain of
equality tests, evaluated top to bottom --- algorithmically a sequence of \texttt{if}
statements, not a single dispatch (\texttt{visit\_Switch}, ir\_gen.py:776--810):
\begin{lstlisting}[style=pseudostyle]
function LOWER-SWITCH(value, cases, default, body):
    end_lbl <- new label "sw_end";  push break->end_lbl
    for each (caseValue, label) in cases:           // in source order
        diff <- value - caseValue
        if diff == 0:  goto label                    // IRCondJump per case
    if default exists:  JUMP default.label
    else:                JUMP end_lbl
    for each (label, stmts) in cases:  LABEL label;  LOWER(stmts)
    if default exists:  LABEL default.label; LOWER(default.stmts)
    LABEL end_lbl
    pop break target
\end{lstlisting}
Because there is no dispatch table, a \texttt{switch} with $N$ cases costs $N$ sequential
comparisons in the worst case (the last \texttt{case} matched) --- the same asymptotic cost
as an equivalent \texttt{if}/\texttt{else if} chain, which is exactly what it compiles to.

\subsection*{\texttt{break} / \texttt{continue}}
Trivial once the enclosing construct has pushed its targets: each is a single
\texttt{IRJump} to whichever label is currently on top of a one-deep
\texttt{\_break\_to}/\texttt{\_cont\_to} pair of fields (ir\_gen.py:820--824), saved and
restored around every loop/switch so nested constructs target the \emph{innermost}
enclosing one:
\begin{lstlisting}[style=pseudostyle]
function LOWER-BREAK():     if break_target exists:  JUMP break_target
function LOWER-CONTINUE():  if continue_target exists:  JUMP continue_target
\end{lstlisting}

\subsection*{Label-naming convention, at a glance}
\begin{table}[h]
\centering
\caption{Label prefixes by construct (all generated by \texttt{ir\_gen.py}'s \texttt{\_lbl()})}
\small
\begin{tabular}{@{}p{2.2cm}p{9.6cm}@{}}
\toprule
Construct & Labels generated \\
\midrule
\texttt{if} & \texttt{if\_t} (true body), \texttt{if\_f} (else body, if present), \texttt{if\_e} (join point) \\
\texttt{while} & \texttt{wc} (condition check), \texttt{wb} (body), \texttt{we} (exit) \\
\texttt{do-while} & \texttt{dw\_b} (body), \texttt{dw\_c} (condition check), \texttt{dw\_e} (exit) \\
\texttt{for} & \texttt{fc} (condition), \texttt{fb} (body), \texttt{fi} (increment), \texttt{fe} (exit) \\
\texttt{switch} & \texttt{case} (one per \texttt{case} label), \texttt{dflt} (default), \texttt{sw\_end} (exit) \\
\texttt{\&\&}/\texttt{||} & \texttt{andF}/\texttt{andE}, \texttt{orT}/\texttt{orE} (short-circuit result materialization) \\
bare comparison & \texttt{cT}/\texttt{cE} (only when a comparison is used as an ordinary 0/1-valued expression, not as a direct branch condition) \\
\bottomrule
\end{tabular}
\end{table}

% ============================================================================
\section{Condition Lowering: From a C Expression to a Single Branch}
% ============================================================================

Every construct above calls one shared helper, \texttt{\_emit\_cond} (ir\_gen.py:1277--1284),
to turn a C boolean expression into IR jumps. Its strategy has two paths:

\begin{lstlisting}[style=pseudostyle, caption={\texttt{\_emit\_cond}, ir\_gen.py:1277--1284}]
function EMIT-COND(cond, true_lbl, false_lbl):
    if cond is a direct comparison (>, <, >=, <=, ==, !=):
        l <- EVAL(cond.left);  r <- EVAL(cond.right)
        emit IRCondJump(l, cond.op, r, true_lbl, false_lbl)        // fast path
        return
    v <- EVAL(cond)                                                 // general expression
    emit IRCondJump(v, '!=', 0, true_lbl, false_lbl)                 // "is it truthy?"
\end{lstlisting}

The \textbf{fast path} matters because APARA's actual hardware branch instruction,
\texttt{\$goto}, only ever tests \emph{one register against zero} (with one of six
relational operators) --- it has no two-register-compare form. So whenever a condition is
already a direct comparison (the overwhelmingly common case --- every loop and \texttt{if}
in this report's examples), \texttt{\_emit\_cond} skips materializing a 0/1 boolean value
entirely and emits the comparison directly as an \texttt{IRCondJump}. Only when a condition
is some other kind of expression (a bare variable, a function call result, a pre-computed
\texttt{\&\&}/\texttt{||}) does it fall back to "evaluate it, then test the result against
zero."

\texttt{codegen.py}'s \texttt{\_emit\_cond\_branch} (codegen.py:463--505) then has to bridge
the gap between the IR's general two-operand comparison and the hardware's
one-register-against-zero branch:
\begin{lstlisting}[style=pseudostyle, caption={\texttt{\_emit\_cond\_branch}, codegen.py:463--505 (essence)}]
function EMIT-BRANCH(left, op, right, true_lbl, false_lbl):
    if left and right are both compile-time constants:
        fold the comparison now; emit one unconditional jump; return    // constant folding
    if right == 0:
        if op in {<, <=}:  scratch <- 0 - left;  branch on (scratch, flip(op))
        else:               branch directly on (left, op)
    else:
        scratch <- (left - right)   if op in {>,>=,==,!=}
                    (right - left)   if op in {<,<=}     // sign flip trick
        branch on (scratch, op-or-flipped-op)
    if false_lbl exists:  emit unconditional JUMP false_lbl
\end{lstlisting}
Every two-operand comparison is reduced to a \emph{subtraction} followed by a single
register-vs-zero branch --- the same "synthesize the missing operation from a cheaper
primitive" strategy already seen elsewhere in this compiler (e.g.\ \texttt{\%} synthesized
as \texttt{a - (a/b)*b}). An \texttt{IRCondJump} with no false label at all (the common
case inside a loop, where falling through to the next label \emph{is} the false branch)
skips the trailing unconditional jump --- but every \texttt{IRCondJump} that specifies
\emph{both} a true and a false label, as every \texttt{visit\_While}/\texttt{visit\_For}
condition does, always compiles to exactly \textbf{two} consecutive control-transfer
instructions: one conditional, one unconditional.

% ============================================================================
\section{A Fully Worked, Hardware-Verified Example}
\label{sec:worked}
% ============================================================================

\begin{lstlisting}[language=C, caption={Source}]
int main() {
    int i;
    i = 0;
    while (i < 16) {
        i = i + 1;
    }
    return i;
}
\end{lstlisting}

\begin{lstlisting}[caption={Generated IR (the loop only) --- exactly the \texttt{wc/wb/we} pattern from Section~2}]
wc_1:
  _t2 = &stack[FP-8]
  _t3 = *(_t2+0)
  if _t3 < 16 goto wb_2          (* IRCondJump: true->wb_2, false->we_3, fast path *)
wb_2:
  _t4 = &stack[FP-8]
  _t5 = *(_t4+0)
  _t6 = _t5 + 1
  _t7 = &stack[FP-8]
  *(_t7+0) = _t6
  goto wc_1
we_3:
\end{lstlisting}

\begin{lstlisting}[caption={Resulting bundled mcode (loop only) --- 9 bundles, hardware-verified}]
wc_1:
||
    + $r2 ($i64) $r26 -8          // r2 = &i
;
||
    $ld ($i32) $r3 [$r2 + 0]      // r3 = i
    + $r4 ($i64) $r0 16           // r4 = 16
;
||
    - $r5 ($i64) $r4 $r3          // r5 = 16 - i  (sign-flip trick for '<')
;
||
    ? ($i64) $r5 > $goto wb_2     // branch TRUE: 16-i > 0  <=>  i < 16
;
||
    ? ($i64) $r0 == $goto we_3    // branch FALSE: unconditional fall-through to exit
;
wb_2:
||
    + $r4 ($i64) $r26 -8
;
||
    $ld ($i32) $r5 [$r4 + 0]
;
||
    + $r6 ($i64) $r5 1            // r6 = i + 1
    + $r7 ($i64) $r26 -8
;
||
    $st ($i32) [$r7 + 0] $r6      // store i+1 back
    ? ($i64) $r0 == $goto wc_1    // loop-back -- SHARES a bundle with the store
;
we_3:
\end{lstlisting}

Two details worth calling out explicitly, both confirmed directly from this listing:
\begin{itemize}
  \item The condition \texttt{i < 16} becomes a \emph{subtraction} (\texttt{\$r5 = \$r4 -
        \$r3}, i.e.\ \texttt{16 - i}) followed by a branch on \texttt{\$r5 > 0} --- exactly
        the sign-flip strategy from \texttt{EMIT-BRANCH} above, because APARA's
        \texttt{\$goto} cannot compare two registers directly.
  \item The loop-back jump (\texttt{? \$r0 == \$goto wc\_1}) does \emph{not} need its own
        bundle: it shares one with the preceding \texttt{\$st}, because the bundler's only
        rule is that a control-transfer instruction must be \textbf{last} in its bundle, not
        that it must be \textbf{alone} in it (Section~\ref{sec:mcode}).
\end{itemize}

% ============================================================================
\section{mcode-Level Handling: One Instruction, Two Hard Rules}
\label{sec:mcode}
% ============================================================================

By the time control flow reaches mcode, every branch in this report --- regardless of
which C construct produced it --- is one of exactly two instruction shapes:
\begin{itemize}
  \item \textbf{Conditional}: \texttt{? (\$iN) \$reg OP \$goto label} --- six possible
        \texttt{OP}s (\texttt{>}, \texttt{<}, \texttt{>=}, \texttt{<=}, \texttt{==}, \texttt{!=}),
        always comparing one register against an implicit zero.
  \item \textbf{Unconditional}: \texttt{? (\$i64) \$r0 == \$goto label} --- there is no
        dedicated unconditional-jump instruction on this ISA, so the compiler reuses the
        conditional form with the always-true tautology \texttt{0 == 0}, needing no extra
        register at all (\texttt{\_emit\_jump}, codegen.py:459--461).
\end{itemize}

\texttt{bundler.py} enforces two rules specific to these instructions (Section
2.4 of the main project report covers the full hazard-detection algorithm; only the
control-flow-specific pair is repeated here):
\begin{enumerate}
  \item \textbf{A control-transfer instruction must be the last instruction in its
        bundle.} Once one has been added to the bundle being packed, the very next
        instruction --- whatever it is --- forces a new bundle to start. It does \emph{not}
        need to be the \emph{only} instruction, as Section~\ref{sec:worked}'s loop-back
        example shows directly.
  \item \textbf{A label forces a new bundle boundary.} Every label (\texttt{wc\_1},
        \texttt{wb\_2}, \texttt{we\_3}, \ldots) attaches to the start of whatever bundle
        comes right after it; two different labels can collapse onto the very same bundle
        only if nothing meaningful separates them, in which case
        \texttt{\_merge\_duplicate\_labels} rewrites every reference to the dropped label
        so it points at the surviving one (this is the exact mechanism --- and the exact
        bug, before it was fixed --- behind the duplicate-label assembler crash documented
        in the project's main bug report).
\end{enumerate}

% ============================================================================
\section{Summary}
% ============================================================================

Every C control-flow construct --- however different it looks in source --- is reduced by
\texttt{ir\_gen.py} to the same three IR primitives (\texttt{IRLabel}, \texttt{IRJump},
\texttt{IRCondJump}), arranged in a construct-specific but fixed pattern of labels.
\texttt{codegen.py} then reduces every two-operand comparison to a subtraction plus a
single register-against-zero test, because that is the only branch form APARA's hardware
actually supports, and reduces unconditional jumps to a tautological comparison, because
APARA has no dedicated unconditional-jump instruction at all. The bundler's only
control-flow-specific obligations are to keep each control-transfer last in its bundle and
to treat every label as a hard bundle boundary --- everything else (how many bundles a loop
costs, what shares a bundle with what) falls out of the same general hazard-packing
algorithm used for arithmetic and memory instructions throughout this compiler.

\end{document}
